\chapter{Critical Analysis, Limitations and Deviations}
\label{ch:limitations}

This chapter distinguishes the project's deliberate deviations from the limitations that
still require experimental validation. The current state remains promising: the software
architecture is consistent, the main nodes exist, and the Kinect/Nav2/\gls{ekf} choices
provide a technically sound foundation. The points below mostly concern the fine
calibration needed to turn a functional integration into a repeatably demonstrable
product.

\section{Architectural Trade-offs}
\label{sec:lim_archi}

During development, several adjustments were made relative to the initial concepts:

\begin{table}[H]
	\centering
	\caption{Summary of architectural changes.}
	\label{tab:ecarts_archi}
	\renewcommand{\arraystretch}{1.3}
	\begin{tabularx}{\linewidth}{L{3.8cm} L{8.4cm}}
		\toprule
		\textbf{Aspect} & \textbf{Final implementation} \\
		\midrule
		Motor/servo control & Centralized on the Raspberry\,Pi\,4; the ESP32-CAM is dedicated to the video stream.\\
		Mapping sensor & Kinect v1 via \texttt{kinect\_ros2} and \texttt{depthimage\_to\_laserscan}, with \texttt{depth\_cleanup} ready for a cleaned stream.\\
		Ultrasonic radar & Kept in the code for testing or as a slow fallback, but not retained as the primary autonomous driving sensor.\\
		Exploration & \texttt{m-explore-ros2} referenced by the branch, still to be orchestrated in a complete EsiBot launch.\\
		Web gateway & A dedicated Python HTTP server for telemetry, simpler than the WebSocket gateway originally planned.\\
		\bottomrule
	\end{tabularx}
\end{table}

These deviations do not degrade the product's coherence. On the contrary, they clarify a
more realistic architecture: real-time control stays on the Raspberry\,Pi, while heavy
perception and planning remain offloadable to the operator PC.

\section{Cross Odometry--SLAM Calibration}
\label{sec:lim_calibration}

The deepest limitation is not a missing module, but the numerical agreement between the
topics that describe the same motion. A correct PWM$\to$RPM calibration is not enough if
the wheel radius, wheelbase, yaw correction, \gls{ekf} fusion and \texttt{slam\_toolbox}
thresholds do not converge on the same real-world displacement.

\begin{table}[H]
	\centering
	\caption{Parameters to be calibrated jointly.}
	\label{tab:calibration_cross_topics}
	\begin{tabularx}{\linewidth}{L{3.4cm} L{4.3cm} L{4.5cm}}
		\toprule
		\textbf{Chain} & \textbf{Sensitive parameters} & \textbf{Symptom if miscalibrated} \\
		\midrule
		Wheels $\to$ \topic{/odom} & calibrator radius 0.033\,m, \texttt{wheel\_base}=0.16\,m, effective CPR & biased integrated distance or rotation \\
		\gls{urdf}/\gls{tf} & 0.025\,m radius in the model, visual wheelbase 0.14\,m & inconsistency between the \gls{rviz} model, costmaps and odometry \\
		\gls{imu} $\to$ \gls{ekf} & gyroscope bias, yaw thresholds, \texttt{imu0\_config} & robot over- or under-corrected when driving straight \\
		\topic{/scan} $\to$ map & \texttt{minimum\_travel\_distance}, \texttt{minimum\_travel\_heading}, Kinect range & stretched map, duplicated walls or unstable loop closures \\
		Nav2 $\to$ \topic{/cmd\_vel} & \gls{dwb} limits, min. speed 0.2\,m/s, accelerations & oscillations or trajectories too aggressive for the motors \\
		\bottomrule
	\end{tabularx}
\end{table}

The recommended method is manual and based on short \ros{} bags: record
\topic{/cmd\_vel}, \topic{/odom}, \topic{/odometry/filtered}, \topic{/scan}, \topic{/tf}
and \topic{/map} over a straight line, an in-place rotation, and a reference square;
compare the actual distance, the odometric trajectory, and the
\topic{map}$\to$\topic{odom} correction; then adjust the radius, wheelbase, and yaw gains
before modifying the \texttt{slam\_toolbox} thresholds. This step requires human
analysis, since the same visible defect on the map can stem from a poorly calibrated
wheel, a biased \gls{imu} yaw, a noisy Kinect scan, or an excessively high update
threshold.

\section{Perception Sensors}
\label{sec:lim_sensors}

The switch from the ultrasonic radar to the Kinect greatly improves autonomous driving,
but does not turn the system into a LiDAR. The remaining limitations are known and
calibratable:

\begin{itemize}
	\item \textbf{Ultrasonic radar} --- the current sweep is slow ($\approx$0.07\,Hz),
	      narrow ($60^\circ$) and sequential. It does not provide dense enough
	      observation for a dynamic costmap; it should remain a test tool or a very
	      local fallback.
	\item \textbf{Kinect} --- the horizontal field of view of about $57^\circ$, the
	      usable 0.35--4.0\,m range and sensitivity to infrared light require active
	      rotation of the robot and indoor use. These limitations are acceptable for the
	      demonstration scenario, since the Kinect still provides a depth scene far more
	      usable than a single ultrasonic beam.
	\item \textbf{\texttt{depth\_cleanup}} --- the library exists and its tests protect
	      the metric format, but the adapter \ros{} node has not yet been inserted into
	      \texttt{mapping.launch.py}. The gain on \topic{/scan} must therefore be
	      measured before the defense.
\end{itemize}

\section{Safety and Operations}
\label{sec:lim_safety}

\begin{itemize}
	\item \textbf{Teleoperation without an enable button} --- the Xbox configuration
	      accepts a command as soon as the stick moves. This is convenient during
	      testing, but a demonstration mode should reintroduce a dead-man button or a
	      clear software stop.
	\item \textbf{No battery supervision} --- no node publishes
	      \topic{/battery\_state}. Long trials must therefore be monitored outside the
	      software to avoid a cutout during navigation.
	\item \textbf{Hardcoded ESP32-CAM WiFi credentials} --- acceptable for a local
	      prototype, but should be replaced with a configuration mechanism if the
	      repository becomes public or is reused.
\end{itemize}

\section{Performance and Orchestration}
\label{sec:lim_perf}

Simultaneously running the Kinect, \texttt{depthimage\_to\_laserscan},
\texttt{slam\_toolbox}, Nav2 and exploration may exceed the comfortable capacity of a
Raspberry\,Pi\,4 alone. The three-layer separation therefore remains relevant: the
Raspberry\,Pi drives the hardware and publishes \topic{/odom}/\topic{/imu/data\_raw},
while the operator PC can carry the Kinect, \gls{slam}, Nav2, \gls{rviz} and
exploration. The point to validate is not only CPU load, but also DDS network latency
and \gls{tf} freshness, especially when exploration sends goals while the map is being
built.

\section{Summary}
\label{sec:lim_synthese}

The product remains optimistic in its expected results: the chosen architecture is
better than the initial ultrasonic plan for credible autonomous navigation. The main
condition is to treat calibration as a complete experimental campaign, not as an
isolated adjustment of the left or right motor. Once this campaign is completed, the
remaining limitations mainly concern operational safety and the orchestration of the
high-level modules.
